\section{Introduction}
Reusability of ontologies depends heavily on documentation and on the availability of representative instance data~\cite{alharbi2021assessing}. However, in emerging or privacy-sensitive domains, such data is often missing, incomplete or difficult to share, which limits empirical ontology evaluation and pushes validation back toward expert judgement~\cite{brank2005survey}. Tool support is already central in ontology engineering~\cite{10.3233/SW-190382}, and recent LLM-based work has mostly focused on ontology design assistance, reasoning support or taxonomy construction~\cite{damiano2025evaluating,yang2026large,soares2025exploring}.

OWL-SDA instead targets ontology evaluation and documentation. It uses AI agents with tools and structured tasks~\cite{yehudai2025survey,10.1145/3491101.3519665,WEBBER1995253} to generate RDF examples when real data cannot be shared because of sensitivity, governance or context-window constraints~\cite{YAN2025100300,wang2024corag,vallevik2025processing,Kaabachi2025,Naik_2023}, or when creating data examples would be too costly or time-consuming. Rather than proposing a new synthetic-data model, OWL-SDA combines shape-level generation, iterative repair and role separation for ontologies expressed in OWL~\cite{Antoniou2009}. The approach fits within LLM-based synthetic data generation~\cite{electronics13173509}, but prioritises structurally valid and interpretable linked-data examples while addressing incompleteness and inaccuracy risks~\cite{hao2024synthetic}. 

We position OWL-SDA against existing work in Section~\ref{sec:related}, and detail our workflow in Section~\ref{sec:solution}. In Section~\ref{sec:evaluation}, we evaluate OWL-SDA against two contrasting scenarios.